% Generated draft: requires mathematical and source review.
\documentclass[11pt]{article}
\usepackage[margin=25mm]{geometry}
\usepackage{fontspec}
\setmainfont{Latin Modern Roman}
\usepackage{amsmath,amssymb,mathtools,unicode-math}
\setmathfont{Latin Modern Math}
\usepackage{enumitem,needspace,xurl,hyperref}
\setlength{\parindent}{0pt}
\setlength{\parskip}{5pt}
\setlength{\emergencystretch}{3em}
\newcommand{\N}{\mathbb{N}}
\newcommand{\Z}{\mathbb{Z}}
\newcommand{\Q}{\mathbb{Q}}
\newcommand{\R}{\mathbb{R}}
\newcommand{\C}{\mathbb{C}}
\newcommand{\F}{\mathbb{F}}
\newcommand{\A}{\mathbb{A}}
\newcommand{\PP}{\mathbb{P}}
\newcommand{\E}{\mathbb{E}}
\newcommand{\eps}{\varepsilon}
\newcommand{\dd}{\,\mathrm{d}}
\begin{document}
\section*{UnsolvedMath ID: 801162}
\textbf{Source problem code:} catalogue entry
\subsection{Definitions and mathematical statement}
Let \(\mathbb K\) denote the scalar field intended in the source, and let \(H_{\mathbb K}=L^2(\mathbb R;\mathbb K)\), with Lebesgue measure and norm
\[
\|f\|_2=\left(\int_{\mathbb R}|f(x)|^2\,\dd x\right)^{1/2}.
\]
Under the usual dyadic dilation and integer translation convention, an orthonormal wavelet is a function \(\psi\in H_{\mathbb K}\) for which
\[
\left\{\psi_{j,k}:j,k\in\mathbb Z\right\},\qquad
\psi_{j,k}(x)=2^{j/2}\psi(2^j x-k),
\]
is an orthonormal basis of \(H_{\mathbb K}\). Write \(\mathcal W_{\mathbb K}\) for the collection of these wavelets. Every element of \(\mathcal W_{\mathbb K}\) has norm \(1\), so this collection is a subset of the unit sphere of \(H_{\mathbb K}\).

The question attributed in the supplied record to D. Larson and G. Weiss, independently around 1995, is whether this collection is path connected in the \(L^2\) norm. Explicitly, is it true that for every \(\psi_0,\psi_1\in\mathcal W_{\mathbb K}\) there exists a map
\[
\gamma:[0,1]\longrightarrow\mathcal W_{\mathbb K}
\]
such that
\[
\gamma(0)=\psi_0,\qquad \gamma(1)=\psi_1,\qquad
\lim_{\substack{t\to t_0\\t\in[0,1]}}
\|\gamma(t)-\gamma(t_0)\|_2=0
\quad\text{for every }t_0\in[0,1]?
\]
The supplied excerpt does not explicitly define the dilation and translation convention or specify whether \(L^2(\mathbb R)\) has real or complex scalars. The displayed wavelet system makes the usual dyadic convention explicit, subject to source confirmation; the intended choice \(\mathbb K=\mathbb R\) or \(\mathbb K=\mathbb C\) remains unresolved here. No equivalence between these scalar-field versions is asserted.
\subsection{Short English statement}
Can every pair of orthonormal wavelets be joined by a curve of orthonormal wavelets that is continuous in the \(L^2(\mathbb R)\) norm? The displayed formulation uses the usual dyadic wavelet convention; the source excerpt leaves the scalar field unspecified.
\subsection{Sources}
\begin{itemize}
\item Local UnsolvedMath registry, permanent ID \(801162\), exact source problem code \texttt{catalogue entry}, title: \emph{Open Problems in Wavelet Theory --- Problem 1}. The supplied registry labels the problem open; this label is source data, not an independently verified status assessment.
\item Bownik and Rzeszotnik, \emph{Open Problems in Wavelet Theory} (2019), Problem 1. Supplied source URL: \url{https://arxiv.org/abs/1905.08230}. The record reports access on 2026-10-16 and extraction by \texttt{arxiv\_latex}; these are supplied provenance details, not a claim of independent access or verification.
\item The supplied statement attributes the question independently to D. Larson and G. Weiss around 1995. The registry lists Bownik and Rzeszotnik with proposed year 2019, corresponding to the cited source list; the earlier attribution is retained separately.
\end{itemize}
\end{document}
